% ===========================================================================
% Chapitre 4 — Implémentation.
%
% Organisation PAR SOUS-SYSTÈME. Chaque section suit le même gabarit :
% objectif, exigences couvertes, conception retenue, difficulté rencontrée.
% Les difficultés rapportées sont RÉELLES et traçables dans le dépôt.
% ===========================================================================
\chapter{Implémentation}
\label{chap:implementation}

\section{Introduction}
\label{sec:impl_intro}

Ce chapitre décrit la réalisation de l'architecture exposée au chapitre
précédent. Il est organisé par sous-système : chaque section énonce l'objectif
du module, les exigences qu'il couvre, la conception retenue, et — lorsqu'elle
existe — la difficulté significative rencontrée et la manière dont elle a été
traitée.

Ce dernier point mérite une remarque de méthode. Un rapport qui ne présente que
des réussites n'informe pas : les décisions y paraissent évidentes alors
qu'elles ne l'étaient pas. Les difficultés rapportées ici sont celles qui ont
réellement modifié la conception.

\section{Organisation de l'implémentation}
\label{sec:impl_organization}

Le tableau~\ref{tab:impl_layers} récapitule l'organisation du code livré.

\begin{table}[htbp]
  \centering
  \caption{Organisation des couches d'implémentation}
  \label{tab:impl_layers}
  \begin{tabularx}{\textwidth}{L{3.2cm} L{2.2cm} X}
    \toprule
    \textbf{Couche} & \textbf{Volume} & \textbf{Contenu} \\
    \midrule
    Interface client   & —            & Composants autonomes, services d'accès, gardes de route, modèles typés \\
    Endpoints          & 18 modules   & 173 routes et deux points de liaison temps réel \\
    Services           & 33 modules   & Règles métier, un module par domaine fonctionnel \\
    Modèles            & 16 modules   & 25 tables et leurs relations \\
    Migrations         & 19 révisions & Source unique de vérité du schéma \\
    Intégrations       & 3 clients    & Services externes, tous facultatifs \\
    \bottomrule
  \end{tabularx}
\end{table}

\section{Socle applicatif}
\label{sec:impl_foundation}

\paragraph{Objectif et exigences.} Fournir configuration typée, gestion du
schéma et journalisation. Couvre NFR-06 et NFR-09.

\paragraph{Conception retenue.} La configuration est validée au démarrage par
un schéma typé : une variable manquante ou mal formée arrête le service avec
un message explicite, plutôt que de produire une défaillance différée au
premier appel qui en dépend. Le schéma de base est géré exclusivement par
migrations (décision D-01) : au démarrage, l'application compare la révision
portée par la base à celle attendue par le code et journalise une erreur en
cas d'écart.

\paragraph{Difficulté rencontrée.} La journalisation applicative était
silencieuse. Le serveur ne configure que ses propres journaux et laisse le
journal racine à un niveau qui écarte les messages d'information : tous les
événements applicatifs — modèle chargé, courriel émis, entretien synchronisé —
étaient rejetés sans trace. Le diagnostic a été retardé par le fait que rien
n'échouait : les traitements fonctionnaient, ils étaient simplement muets. La
correction configure explicitement le journal racine au démarrage, sans
perturber celui du serveur.

\section{Identité, authentification et second facteur}
\label{sec:impl_authentication}

\paragraph{Objectif et exigences.} Authentifier les utilisateurs et protéger
les comptes. Couvre FR-AUTH-01 à FR-AUTH-05 et NFR-01.

\paragraph{Conception retenue.} Les mots de passe sont hachés par une fonction
lente et salée, paramétrée à douze tours — un compromis entre coût d'attaque et
temps de réponse acceptable. Le second facteur repose sur un mot de passe à
usage unique fondé sur le temps.

\paragraph{Deux décisions issues d'un modèle de menace.} Elles méritent d'être
détaillées, car elles ne relèvent pas de la préférence.

D'abord, l'\emph{égalisation du temps de réponse}. Lorsqu'une adresse inconnue
est soumise, le système calcule malgré tout un hachage factice avant de
répondre. Sans cette précaution, la réponse serait sensiblement plus rapide
pour une adresse inexistante — le hachage étant délibérément lent — ce qui
permettrait d'énumérer les comptes enregistrés en mesurant les délais. Le
contrôle ne protège pas le mot de passe : il protège l'information
« cette adresse existe ».

Ensuite, le \emph{plafonnement des tentatives par compte} et non par adresse
réseau. Un plafond par adresse se contourne en changeant de source, ce qui est
trivial ; un plafond par compte visé résiste à cette rotation. Le compromis
assumé est qu'un attaquant peut ainsi verrouiller le compte d'un tiers : c'est
une gêne, préférable à une prise de contrôle.

\section{Autorisation par rôle}
\label{sec:impl_authorization}

\paragraph{Objectif et exigences.} Garantir qu'une action n'est possible que
pour les rôles habilités. Couvre FR-ACC-01, FR-ACC-02 et NFR-02.

\paragraph{Conception retenue.} La politique est déclarative et externalisée
(décision D-02). Elle exprime des permissions métier et non des chemins
d'interface, et porte un héritage entre rôles : le responsable RH hérite du
recruteur, l'administrateur du responsable RH. Une permission accordée au
recruteur remonte donc automatiquement, et la politique ne répète pas trois
fois la même ligne. Côté client, des gardes de route masquent les entrées de
menu inaccessibles — non comme un contrôle, mais comme un confort : la décision
reste serveur, et un client qui mentirait sur ses droits n'obtiendrait rien.

\paragraph{Écart constaté.} La campagne du 22 septembre 2026 a révélé qu'un
endpoint engageant financièrement l'entreprise est gardé par une permission
\emph{plus large} que celle que la politique entend exiger. L'écart n'est pas
dans le moteur — qui applique exactement ce qui lui est déclaré — mais entre
l'intention de la politique et la déclaration de la route. Ce constat est
analysé à la section~\ref{sec:failure_analysis}. Il illustre la valeur de la
décision D-02 : la politique étant lisible en un seul fichier, l'écart a pu
être localisé par comparaison, ce qu'une logique dispersée en conditions
n'aurait pas permis.

\section{Offres, candidatures et gestion des CV}
\label{sec:impl_recruitment}

\paragraph{Objectif et exigences.} Gérer les offres, les candidatures et les
fichiers déposés. Couvre FR-OFF-01 à FR-OFF-05 et FR-CAND-01 à FR-CAND-07.

\paragraph{Conception retenue.} Le dépôt est contrôlé \emph{avant} tout
traitement : liste blanche de types et plafond de taille. Ces deux contrôles
sont appliqués côté serveur, le client n'étant qu'une commodité. La suppression
d'une candidature déclenche le nettoyage du fichier associé, pour éviter
l'accumulation de documents orphelins — qui seraient, s'agissant de CV, une
exposition durable de données personnelles.

\pqscreenshot{job_list_public}
  {Catalogue public des offres d'emploi avec recherche et filtres.}
  {fig:screen_joblist}

\pqscreenshot{job_detail_apply}
  {Fiche d'une offre et formulaire de candidature avec dépôt de CV.}
  {fig:screen_apply}

\paragraph{Confidentialité.} Aucune donnée extraite d'un CV n'est écrite dans
les journaux, à aucun niveau de verbosité. Cette contrainte a orienté
l'écriture des messages de diagnostic : ils rapportent des identifiants et des
compteurs, jamais des extraits de contenu. Un message de débogage affichant
« compétences extraites : \ldots » aurait été commode au développement et
inacceptable en exploitation.

\section{Analyse de CV et notation des candidats}
\label{sec:impl_scoring}

\paragraph{Objectif et exigences.} Extraire les faits d'un CV et noter la
candidature. Couvre FR-CV-01 à FR-CV-03 et FR-SCORE-01 à FR-SCORE-05.

\paragraph{Conception retenue.} L'analyse s'exécute en tâche de fond
(décision D-08), dans une session de base dédiée. Elle extrait le texte, puis
les compétences, l'expérience et la formation. Le score est ensuite calculé par
comparaison aux exigences de l'offre et \emph{stocké décomposé} : les
sous-scores sont persistés, non recalculés à l'affichage. Cette décision
garantit que ce qui a été montré au recruteur reste consultable ensuite, même
si la logique de notation évolue.

\pqscreenshot{cv_analysis_breakdown}
  {Décomposition du score d'un candidat en sous-scores explicables.}
  {fig:screen_cvanalysis}

\pqscreenshot{recruitment_ranking}
  {Classement des candidats d'une offre selon leur score de correspondance.}
  {fig:screen_ranking}

\paragraph{Difficulté rencontrée : une surnotation silencieuse.} C'est le
défaut le plus instructif du projet.

La normalisation des compétences requises comportait un mécanisme de repli :
lorsqu'une compétence de l'offre n'était pas reconnue dans le référentiel, un
découpage de secours permettait de l'exploiter malgré tout. Ce repli était
appliqué \emph{globalement} — c'est-à-dire déclenché seulement si
\emph{aucune} compétence n'était reconnue.

La conséquence était subtile. Dès qu'une seule compétence courante était
reconnue — « Python », par exemple — le repli ne se déclenchait pas, et toutes
les autres compétences non reconnues, typiquement un cadriciel interne ou un
outil métier, étaient silencieusement perdues. L'offre exigeait donc moins que
ce qu'elle demandait réellement, et les candidats étaient \textbf{surnotés}.

Le défaut ne produisait aucune erreur : les scores étaient plausibles,
simplement trop élevés, et rien ne signalait les compétences disparues. Il a
été mis en évidence en construisant un cas de test où une offre mêlait
délibérément des compétences standard et des compétences inventées, puis en
vérifiant que le nombre de compétences effectivement prises en compte
correspondait au nombre déclaré.

La correction applique le repli \emph{compétence par compétence} plutôt que
globalement. Chaque compétence non reconnue emprunte désormais son propre
chemin de secours, indépendamment du sort des autres.

\section{Rapprochement sémantique}
\label{sec:impl_semantic}

\paragraph{Objectif et exigences.} Retrouver des profils par proximité de sens.
Couvre FR-SEM-01 à FR-SEM-04.

\paragraph{Conception retenue.} Les CV analysés sont segmentés puis convertis
en vecteurs de dimension 384, selon l'approche par réseaux siamois qui rend
les plongements de phrases directement comparables par similarité
cosinus~\cite{reimers2019sbert}. La recherche emprunte deux chemins
(décision D-03) : un index vectoriel natif~\cite{pgvector} lorsque l'extension
de base de données est disponible, et un repli applicatif calculant la
similarité en mémoire sinon. Le service détecte la disponibilité au démarrage et aiguille en
conséquence.

L'indexation est délibérément traitée en \emph{meilleur effort} : un échec
d'indexation ne remet pas en cause la notation (FR-SEM-04). L'analyse reste
exploitable, la recherche sémantique est simplement moins riche. L'inverse —
annuler une notation parce qu'une indexation a échoué — aurait fait dépendre
une fonction essentielle d'une fonction accessoire.

\pqscreenshot{semantic_search_results}
  {Recherche sémantique de profils : requête en langage naturel et résultats
   classés par proximité de sens.}
  {fig:screen_semantic}

\paragraph{Limite de validation à signaler.} La campagne du 22 septembre 2026
n'a exercé que le \emph{chemin de repli}. Le chemin vectoriel natif est
implémenté mais n'a pas été capturé en exécution. Affirmer que la recherche
native fonctionne serait une extrapolation.

\section{Assistant extractif}
\label{sec:impl_cag}

\paragraph{Objectif et exigences.} Répondre aux questions internes à partir
d'une base de connaissances. Couvre FR-CAG-01 et FR-CAG-02.

\paragraph{Conception retenue.} L'assistant est \emph{extractif} et non
génératif (décision D-06). Il reprend la phase de \emph{récupération} du
schéma de génération augmentée par récupération~\cite{lewis2020rag}, mais
s'arrête avant la phase de génération : il indexe la base de connaissances,
retrouve les passages pertinents et les restitue \emph{en citant leur
source}. Il ne reformule pas.

\paragraph{Justification du choix.} Une réponse générée est fluide mais
invérifiable : le lecteur ne peut pas distinguer un fait issu de la base d'une
formulation plausible produite par le modèle. Dans un contexte de recrutement,
où une réponse peut orienter une décision concernant une personne, cette
indistinction est un défaut et non une commodité. Le coût assumé est une
réponse moins naturelle ; le bénéfice est qu'elle est traçable jusqu'à sa
source.

\pqscreenshot{cag_assistant_sources}
  {Assistant extractif : réponse accompagnée des sources dont elle est tirée.}
  {fig:screen_cag}

\section{Entretiens et visioconférence}
\label{sec:impl_interviews}

\paragraph{Objectif et exigences.} Planifier, confirmer et conduire les
entretiens. Couvre FR-INT-01 à FR-INT-08.

\paragraph{Conception retenue.} Le recruteur propose jusqu'à cinq créneaux et
désigne le recruteur qui conduira l'entretien. Le candidat reçoit un lien
public à jeton et confirme un créneau sans compte. À la confirmation, le
système vérifie qu'aucun autre candidat n'a réservé le même créneau entre
temps, avertit le recruteur assigné, génère le lien de visioconférence et
synchronise l'agenda.

Le lien de visioconférence suit un repli : lorsque le service principal ne
délivre pas de salle, une salle de secours est dérivée du jeton d'invitation —
un lien stable par entretien, sans compte ni configuration.

\pqscreenshot{interview_scheduler_recruiter}
  {Planification d'un entretien : créneaux proposés et désignation du
   recruteur assigné.}
  {fig:screen_scheduler}

\pqscreenshot{interview_manager_slots}
  {Gestion des entretiens : cartes par statut, pagination et état de la
   synchronisation d'agenda.}
  {fig:screen_interviews}

\paragraph{Cloisonnement par recruteur.} Un recruteur ne voit que les
entretiens qui lui sont affectés (FR-INT-04). Le filtrage est appliqué
\emph{côté serveur}, à partir du rôle de l'appelant, et non par masquage dans
l'interface : un filtre client serait contournable par un appel direct.

\paragraph{Difficulté rencontrée : une session de base déjà fermée.} La
synchronisation d'agenda, exécutée en tâche de fond après la réponse, échouait
lorsqu'elle tentait de relire l'invitation : la session de la requête était
close. La correction ouvre une session dédiée dans la tâche, selon le même
principe que l'analyse de CV. C'est un schéma récurrent : tout traitement
survivant à la réponse doit posséder sa propre session.

\section{Test technique}
\label{sec:impl_assessment}

\paragraph{Objectif et exigences.} Évaluer les compétences par un test
chronométré. Couvre FR-TEST-01 à FR-TEST-04.

\paragraph{Conception retenue.} Le test est proposé par lien public à jeton. La
session est chronométrée, les réponses corrigées automatiquement, et un
balayage périodique clôt les sessions dont le temps est écoulé — sans quoi un
candidat fermant son navigateur laisserait sa session ouverte indéfiniment.

\paragraph{Point de conception à souligner.} Le seuil d'éligibilité est évalué
\emph{côté serveur} (FR-TEST-03). Un seuil vivant dans le navigateur finit
toujours par diverger de celui du serveur — au premier déploiement partiel, au
premier cache non invalidé — et la divergence est silencieuse. Placer la règle
là où la décision est prise supprime la question.

\pqscreenshot{test_public_session}
  {Passage du test technique : minuteur, progression et question en cours.}
  {fig:screen_test}

\section{Prédiction salariale et négociation}
\label{sec:impl_salary}

\paragraph{Objectif et exigences.} Estimer une fourchette salariale et assister
la négociation. Couvre FR-SAL-01, FR-SAL-02, FR-NEGO-01 et FR-NEGO-02.

Ce sous-système est le plus démonstratif du projet sur le plan de la rigueur
méthodologique, parce que chaque décision y est mesurée.

\paragraph{Audit du jeu de données.} Les salaires du jeu source étaient
exprimés en net, tandis que la prédiction attendue porte sur le brut. La
conversion n'a pas été faite par un coefficient forfaitaire — qui aurait
introduit une erreur systématique croissante avec le salaire — mais par
\emph{recherche dichotomique sur le barème social réel} : pour chaque salaire
net, on recherche le brut dont la conversion redonne ce net. La démarche est
plus coûteuse et exacte sur toute la plage.

Une seule observation a été exclue de l'entraînement. Elle a été
\emph{conservée dans le fichier source avec son motif écrit}, plutôt que
supprimée : une exclusion silencieuse est indistinguable d'une erreur de
collecte, et devient inauditable quelques mois plus tard.

\paragraph{Sélection du modèle.} Huit modèles ont été comparés sur cinq
partitions indépendantes. La sélection s'est faite sur l'erreur de \emph{test},
et non d'entraînement (décision D-04). Ce choix est déterminant : un modèle
suffisamment souple atteint une erreur d'entraînement quasi nulle en mémorisant
ses exemples, et échoue sur toute donnée nouvelle. Seule l'erreur mesurée sur
des données non vues pendant l'apprentissage distingue l'apprentissage de la
mémorisation.

\paragraph{Retrait d'une variable colinéaire.} Une variable d'entrée était
fortement corrélée à une autre. Sa conservation créait un mode de défaillance
identifié : le modèle répartissait arbitrairement le poids entre les deux
variables, rendant les prédictions instables pour des profils où elles
divergeaient. Son retrait (décision D-05) a supprimé l'instabilité sans dégrader
l'erreur.

\paragraph{Deux défauts d'inférence corrigés.} Le premier : la détection des
compétences procédait par recherche de sous-chaîne, si bien que le langage
\texttt{java} était détecté dans \texttt{JavaScript} — deux technologies
distinctes, conduisant à une estimation erronée. Le second : les intitulés de
poste rédigés en français n'étaient jamais rapprochés d'un barème libellé en
anglais, et retombaient systématiquement sur une valeur par défaut. Les deux
défauts produisaient des prédictions plausibles mais fausses, ce qui les rendait
difficiles à repérer sans jeu de test dédié.

\pqscreenshot{salary_prediction_form}
  {Prédiction salariale : variables d'entrée et fourchette estimée.}
  {fig:screen_salary}

\pqscreenshot{negotiation_live_log}
  {Journal de négociation : tours successifs, montants et arguments échangés.}
  {fig:screen_negotiation}

\section{Contrats et signature électronique}
\label{sec:impl_contracts}

\paragraph{Objectif et exigences.} Éditer, signer et concrétiser l'embauche.
Couvre FR-CTR-01 à FR-CTR-06.

\paragraph{Conception retenue.} Le contrat est édité à partir d'une candidature
retenue, puis envoyé par lien public à jeton. À la signature, le système
calcule une \emph{empreinte des termes} — de sorte qu'une modification
ultérieure du contrat serait détectable — émet un certificat et consigne la
piste d'audit : nom du signataire, horodatage, adresse et agent. La fiche
employé est ensuite créée ou rattachée, avec dédoublonnage par adresse
(décision D-07).

\pqscreenshot{contract_sign_certificate}
  {Contrat signé : certificat, empreinte des termes et piste d'audit.}
  {fig:screen_contract}

\paragraph{Limite juridique à énoncer franchement.} La signature implémentée
est auto-hébergée. Elle établit une \emph{preuve d'intention} — le signataire a
consulté les termes et les a validés, ce que l'empreinte et la piste d'audit
documentent — mais elle ne constitue pas une signature électronique qualifiée
au sens d'un règlement de certification. La présenter autrement serait une
affirmation juridique que le projet n'étaye pas.

\section{Ressources humaines internes}
\label{sec:impl_hr}

\paragraph{Objectif et exigences.} Gérer employés, budgets et carrières.
Couvre FR-RH-01 à FR-RH-03.

\paragraph{Conception retenue.} Le répertoire des employés est alimenté par la
signature des contrats, ce qui évite une ressaisie. Les budgets sont tenus par
département avec imputation des dépenses, et les plans de carrière portent
objectifs et jalons.

\paragraph{Point de modélisation à signaler.} La table des employés ne porte
\emph{aucune clé étrangère} vers celle des comptes. Ce n'est pas un oubli :
un employé enregistré n'a pas nécessairement de compte sur la plateforme. Le
rattachement se fait par adresse lorsque les deux existent. La contrepartie
est que l'intégrité de ce lien n'est pas garantie par le moteur de base, ce qui
constitue une dette assumée.

\section{Rapports, décision et notifications}
\label{sec:impl_reporting}

\paragraph{Objectif et exigences.} Restituer l'activité et diffuser les
événements. Couvre FR-REP-01 à FR-REP-03 et FR-NOTIF-01, FR-NOTIF-02.

\paragraph{Conception retenue.} Les synthèses agrègent l'activité de
recrutement et peuvent être archivées sous forme d'instantanés — un rapport
archivé reste lisible même si les données sous-jacentes évoluent. Les
notifications sont diffusées en temps réel sur une liaison authentifiée par
jeton, la liaison temps réel ne portant pas les en-têtes d'une requête
ordinaire.

\section{Interface utilisateur}
\label{sec:impl_frontend}

\paragraph{Objectif.} Fournir les parcours candidat, recruteur et
administration.

\paragraph{Conception retenue.} L'interface est bâtie en composants autonomes,
avec chargement différé des modules lourds afin de contenir la taille du bundle
initial.

\paragraph{Résultat à ne pas embellir.} Le chargement différé n'a \emph{pas}
suffi. Au 22 septembre 2026, le bundle initial mesure 1,25 Mo bruts pour
246,82 ko transférés, au-delà du budget fixé, et le compilateur émet quatre
avertissements de dépassement. L'exigence NFR-08 n'est pas satisfaite. Présenter
le chargement différé comme un succès serait contredit par la mesure, qui
figure au chapitre~\ref{chap:validation}.

\pqscreenshot{my_applications_status}
  {Suivi des candidatures côté candidat, avec les différents statuts du cycle
   de vie.}
  {fig:screen_myapps}

\pqscreenshot{application_score_detail}
  {Restitution de l'analyse de CV au candidat, avec le score de
   correspondance.}
  {fig:screen_scoredetail}

\section{Conteneurisation et mise à disposition}
\label{sec:impl_packaging}

\paragraph{Objectif.} Rendre la plateforme déployable de façon reproductible.

\paragraph{Conception retenue.} Les images sont construites en plusieurs
étapes, s'exécutent sous un utilisateur non privilégié et exposent des sondes
de santé. La variante allégée de la bibliothèque de calcul tensoriel est
employée, ce qui évite environ 2,5 Go d'image pour des capacités dont le projet
n'a pas l'usage.

\paragraph{Difficulté rencontrée.} La sonde de santé échouait alors que le
service répondait. La sonde interrogeait le nom local, que la résolution
traduisait par une adresse de sixième version, tandis que le service n'écoutait
que sur la quatrième. Le symptôme — un conteneur déclaré non sain mais
parfaitement fonctionnel — ne désignait pas sa cause. La sonde force désormais
la version quatre (décision D-08 du chapitre précédent).

\section{Difficultés rencontrées et ajustements}
\label{sec:impl_challenges}

Le tableau~\ref{tab:impl_challenges} récapitule les difficultés significatives
et les réponses apportées.

\begin{table}[htbp]
  \centering
  \caption{Difficultés d'implémentation et réponses d'ingénierie}
  \label{tab:impl_challenges}
  \begin{tabularx}{\textwidth}{L{3.0cm} X L{3.6cm}}
    \toprule
    \textbf{Sous-système} & \textbf{Difficulté} & \textbf{Réponse} \\
    \midrule
    Socle
      & Journalisation applicative silencieuse : les messages étaient rejetés sans trace
      & Configuration explicite du journal racine \\
    Notation
      & Repli de normalisation appliqué globalement : compétences perdues et candidats \textbf{surnotés}
      & Repli appliqué compétence par compétence \\
    Sémantique
      & Écriture de vecteurs hors du cycle applicatif, invisible à la recherche native
      & Indexation replacée dans le cycle, avec traitement en meilleur effort \\
    Entretiens
      & Session de base fermée avant l'exécution de la tâche de fond
      & Ouverture d'une session dédiée dans la tâche \\
    Salaire
      & Détection de compétence par sous-chaîne : \texttt{java} trouvé dans \texttt{JavaScript}
      & Correspondance sur mot entier \\
    Salaire
      & Intitulés français jamais rapprochés d'un barème anglais
      & Table de correspondance bilingue \\
    Salaire
      & Variable colinéaire produisant des prédictions instables
      & Retrait de la variable, erreur inchangée \\
    Conteneurs
      & Sonde de santé résolue en adresse de sixième version, service en quatrième
      & Sonde forcée sur la quatrième version \\
    Interface
      & Bundle initial au-delà du budget malgré le chargement différé
      & \textbf{Non résolu} — consigné comme limite \\
    \bottomrule
  \end{tabularx}
\end{table}

\section{Traçabilité d'implémentation}
\label{sec:impl_traceability}

Le tableau~\ref{tab:impl_traceability} relie chaque groupe d'exigences à son
sous-système de réalisation. La matrice détaillée figure en
annexe~\ref{app:traceability}.

\begin{table}[htbp]
  \centering
  \caption{Correspondance exigences — implémentation}
  \label{tab:impl_traceability}
  \begin{tabularx}{\textwidth}{L{3.6cm} L{3.4cm} X}
    \toprule
    \textbf{Exigences} & \textbf{Sous-système} & \textbf{Section} \\
    \midrule
    FR-AUTH-01…05        & Authentification        & \ref{sec:impl_authentication} \\
    FR-ACC-01…05         & Autorisation et comptes & \ref{sec:impl_authorization} \\
    FR-OFF-01…05, FR-CAND-01…07 & Offres et candidatures & \ref{sec:impl_recruitment} \\
    FR-CV-01…03, FR-SCORE-01…05 & Analyse et notation & \ref{sec:impl_scoring} \\
    FR-SEM-01…04         & Rapprochement sémantique& \ref{sec:impl_semantic} \\
    FR-CAG-01…02         & Assistant extractif     & \ref{sec:impl_cag} \\
    FR-INT-01…08         & Entretiens              & \ref{sec:impl_interviews} \\
    FR-TEST-01…04        & Test technique          & \ref{sec:impl_assessment} \\
    FR-SAL-01…02, FR-NEGO-01…02 & Salaire et négociation & \ref{sec:impl_salary} \\
    FR-CTR-01…06         & Contrats et signature   & \ref{sec:impl_contracts} \\
    FR-RH-01…03          & RH interne              & \ref{sec:impl_hr} \\
    FR-REP-01…03, FR-NOTIF-01…02 & Rapports et notifications & \ref{sec:impl_reporting} \\
    \bottomrule
  \end{tabularx}
\end{table}

\section{Conclusion du chapitre}
\label{sec:impl_conclusion}

Ce chapitre a décrit la réalisation des treize sous-systèmes de PIQBIT et
rapporté neuf difficultés significatives avec leur traitement. Deux d'entre
elles méritaient un développement : la surnotation silencieuse due à un repli
de normalisation mal placé, et la méthodologie de sélection du modèle
salarial — audit du jeu de données, sélection sur l'erreur de test, retrait
d'une variable colinéaire. Une difficulté reste \emph{non résolue} et est
consignée comme telle : le dépassement du budget de taille du bundle.

Le chapitre suivant présente la validation de la plateforme : stratégie de
test, résultats datés du 22 septembre 2026 — y compris les trois échecs —
qualité du modèle salarial, validation de sécurité et limites d'exploitation.
